\documentclass[aspectratio=169,11pt]{beamer}

% ---------------------------------
% PACKAGES
% ---------------------------------
\usepackage[utf8]{inputenc}
\usepackage[T1]{fontenc}
\usepackage{amsmath}
\usepackage{graphicx}
\usepackage{booktabs}

% ---------------------------------
% THEME AND APPEARANCE
% ---------------------------------
\usetheme{Madrid}
\usecolortheme{default}

\setbeamercolor{structure}{fg=blue!70!black}
\setbeamercolor{frametitle}{bg=blue!70!black,fg=white}
\setbeamercolor{title}{fg=white}
\setbeamercolor{titlelike}{bg=blue!70!black,fg=white}

% ---------------------------------
% PRESENTATION INFORMATION
% ---------------------------------
\title[Beamer Lab]{Introduction to Beamer}

\subtitle{Technical Writing Laboratory}

\author[Ananya Siju]{Ananya Siju}

\institute[MNIT Jaipur]{
Department of Computer Science and Engineering\\
MNIT Jaipur
}

\date[Academic Year 2026--2027]{
Academic Year 2026--2027
}

% ---------------------------------
% BEGIN DOCUMENT
% ---------------------------------
\begin{document}

% =================================
% SLIDE 1 - TITLE PAGE
% =================================
\begin{frame}
    \titlepage
\end{frame}

% =================================
% SLIDE 2 - OUTLINE
% =================================
\begin{frame}{Presentation Outline}

    \tableofcontents

\end{frame}

% =================================
% SECTION 1 - INTRODUCTION
% =================================
\section{Introduction}

% SLIDE 3
\begin{frame}{What is Beamer?}

    \begin{itemize}
        \item Beamer is a LaTeX document class
              for creating slides.

        \item Slides are called frames.

        \item A source file is compiled
              to a portable PDF.

        \item Ideal for mathematical
              and technical presentations.
    \end{itemize}

\end{frame}

% =================================
% SECTION 2 - CORE FEATURES
% =================================
\section{Core Features}

% SLIDE 4
\begin{frame}{Frames and Structured Lists}

    \begin{enumerate}
        \item Declare a frame with
              \texttt{\textbackslash begin\{frame\}}.

        \item Add a title and content.

        \item Close with
              \texttt{\textbackslash end\{frame\}}.
    \end{enumerate}

\end{frame}

% =================================
% SLIDE 5 - BLOCKS
% =================================
\begin{frame}{Blocks and Emphasis}

    \begin{block}{Standard Block}
        Group related ideas in a named block.
    \end{block}

    \begin{alertblock}{Important}
        Keep slides concise and readable.
    \end{alertblock}

    \begin{exampleblock}{Example}
        Use \texttt{\textbackslash textbf}
        for bold text.
    \end{exampleblock}

\end{frame}

% =================================
% SLIDE 6 - OVERLAYS
% =================================
\begin{frame}{Overlays: Step-by-Step Reveal}

    \begin{itemize}

        \item<1->
        First idea appears on overlay 1.

        \item<2->
        Second idea appears on overlay 2.

        \item<3->
        Third idea appears on overlay 3.

    \end{itemize}

\end{frame}

% =================================
% SECTION 3 - TECHNICAL CONTENT
% =================================
\section{Technical Content}

% =================================
% SLIDE 7 - MATHEMATICS
% =================================
\begin{frame}{Mathematical Typesetting}

    The quadratic formula is

    \[
        x =
        \frac{-b \pm \sqrt{b^2-4ac}}{2a},
        \qquad a \neq 0.
    \]

    \vspace{0.5cm}

    An inline expression:
    $E = mc^2$.

\end{frame}

% =================================
% SLIDE 8 - TABLES
% =================================
\begin{frame}{Tables in Beamer}

    \begin{table}

        \centering

        \begin{tabular}{ll}

            \toprule

            Feature & Purpose \\

            \midrule

            Frames & Organize slides \\

            Themes & Control presentation appearance \\

            Overlays & Reveal content progressively \\

            Blocks & Highlight related information \\

            \bottomrule

        \end{tabular}

        \caption{Common Beamer features}

    \end{table}

\end{frame}

% =================================
% SLIDE 9 - COLUMNS
% =================================
\begin{frame}{Two-Column Layout}

    \begin{columns}[T]

        % LEFT COLUMN
        \begin{column}{0.48\textwidth}

            \textbf{Advantages}

            \begin{itemize}
                \item Consistent design
                \item Excellent equations
                \item Version-control friendly
            \end{itemize}

        \end{column}

        % RIGHT COLUMN
        \begin{column}{0.48\textwidth}

            \textbf{Considerations}

            \begin{itemize}
                \item Requires LaTeX syntax
                \item Compilation needed
                \item Fine-grained visual
                      design can take time
            \end{itemize}

        \end{column}

    \end{columns}

\end{frame}

% =================================
% SECTION 4 - WORKFLOW
% =================================
\section{Workflow}

% =================================
% SLIDE 10 - PRACTICAL WORKFLOW
% =================================
\begin{frame}{Practical Workflow}

    \begin{enumerate}

        \item Write a \texttt{.tex} source file.

        \item Select a theme and add frames.

        \item Compile using
              \texttt{pdflatex} or Overleaf.

        \item Review the PDF and fix any errors.

        \item Commit both source and final
              output to GitHub.

    \end{enumerate}

\end{frame}

% =================================
% SLIDE 11 - LEARNING OUTCOMES
% =================================
\begin{frame}{Learning Outcomes}

    \begin{itemize}

        \item Create slides using
              the Beamer class.

        \item Use themes, blocks, columns,
              equations, and tables.

        \item Understand overlay specifications.

        \item Compile a reproducible
              PDF presentation.

    \end{itemize}

\end{frame}

% =================================
% SLIDE 12 - THANK YOU
% =================================
\begin{frame}

    \centering

    \Huge Thank You!

    \vspace{0.6cm}

    \Large Technical Writing Laboratory

    \vspace{0.3cm}

    \large MNIT Jaipur

\end{frame}

% ---------------------------------
% END DOCUMENT
% ---------------------------------
\end{document}

